\diarytitle{0065}{ラプラス方程式の境界値問題（矩形領域、上辺条件が二次関数）}

ラプラス方程式 $u_{xx} + u_{yy} = 0$ は、直交座標では $u = X(x)Y(y)$ の積の形を仮定すると常微分方程式に帰着でき、変数分離法で解ける。

\textbf{手順1: 変数分離と固有値問題。}\ $u = X(x)Y(y)$ とおいて代入し $XY$ で割ると
\[
\frac{X''}{X} = -\frac{Y''}{Y} = -\lambda
\quad\Longrightarrow\quad
X'' + \lambda X = 0, \qquad Y'' - \lambda Y = 0
\]
$u(0,y) = u(\pi,y) = 0$ より $X(0) = X(\pi) = 0$。この固有値問題の解は
\[
\lambda_n = n^2, \qquad X_n(x) = \sin(nx) \qquad (n = 1, 2, 3, \dots)
\]

\textbf{手順2: $Y$ 方向。}\ $Y_n'' - n^2 Y_n = 0$ の一般解を $Y_n = A_n\sinh(ny) + B_n\cosh(ny)$ ととる。$u(x,0) = 0$ より $Y_n(0) = B_n = 0$。よって $Y_n = A_n\sinh(ny)$ であり、重ね合わせて
\[
u(x,y) = \sum_{n=1}^{\infty} A_n \sinh(ny)\sin(nx)
\]

\textbf{手順3: 残った境界条件をフーリエ正弦級数で処理。}\ $y=1$ を代入すると
\[
u(x,1) = \sum_{n=1}^{\infty} \bigl(A_n \sinh n\bigr)\sin(nx) = x(\pi - x)
\]
これは $x(\pi-x)$ を $[0,\pi]$ でフーリエ正弦級数に展開した形なので、係数公式より
\[
A_n \sinh n = \frac{2}{\pi}\int_0^{\pi} x(\pi - x)\sin(nx)\,dx
\]

\textbf{手順4: 積分の計算（部分積分）。}\ まず
\[
\int_0^\pi x\sin(nx)\,dx
= \left[-\frac{x\cos(nx)}{n}\right]_0^\pi + \frac{1}{n}\int_0^\pi \cos(nx)\,dx
= -\frac{\pi(-1)^n}{n} + \left[\frac{\sin(nx)}{n^2}\right]_0^\pi
= -\frac{\pi(-1)^n}{n}
\]
（$\cos(n\pi) = (-1)^n$, $\sin(n\pi) = 0$ を用いた）。次に部分積分を2回行って
\begin{align*}
\int_0^\pi x^2\sin(nx)\,dx
&= \left[-\frac{x^2\cos(nx)}{n}\right]_0^\pi + \frac{2}{n}\int_0^\pi x\cos(nx)\,dx \\
&= -\frac{\pi^2(-1)^n}{n} + \frac{2}{n}\left(\left[\frac{x\sin(nx)}{n}\right]_0^\pi - \frac{1}{n}\int_0^\pi \sin(nx)\,dx\right) \\
&= -\frac{\pi^2(-1)^n}{n} + \frac{2}{n}\left(0 + \left[\frac{\cos(nx)}{n^2}\right]_0^\pi\right)
= -\frac{\pi^2(-1)^n}{n} + \frac{2(-1)^n - 2}{n^3}
\end{align*}
これらより
\[
\int_0^\pi x(\pi-x)\sin(nx)\,dx
= \pi\left(-\frac{\pi(-1)^n}{n}\right) - \left(-\frac{\pi^2(-1)^n}{n} + \frac{2(-1)^n-2}{n^3}\right)
= \frac{2\bigl(1-(-1)^n\bigr)}{n^3}
\]
$n$ が偶数のときは $(-1)^n = 1$ なので $1 - (-1)^n = 0$、$n = 2k+1$（奇数）のときは $(-1)^n = -1$ なので $1 - (-1)^n = 2$ となり、積分値は奇数 $n$ に対してのみ $\dfrac{4}{n^3}$ が残る。したがって
\[
A_n \sinh n = \frac{2}{\pi}\cdot\frac{4}{n^3} = \frac{8}{\pi n^3} \quad (n\text{: 奇数}), \qquad A_n = 0\ (n\text{: 偶数})
\]
よって
\[
\boxed{\;
u(x,y) = \frac{8}{\pi}\sum_{k=0}^{\infty} \frac{\sinh\bigl((2k+1)y\bigr)}{(2k+1)^{3}\sinh(2k+1)}\,\sin\bigl((2k+1)x\bigr)
\;}
\]

\textbf{検算。}\ 各項 $\sin(nx)\sinh(ny)$ は $u_{xx}+u_{yy} = (-n^2+n^2)\sin(nx)\sinh(ny) = 0$ を満たし調和である。$y=0$ で $\sinh 0=0$ より $u(x,0)=0$。$x=0,\pi$ で $\sin(nx)=0$ より $u(0,y)=u(\pi,y)=0$。$y=1$ では $\sinh n/\sinh n = 1$ となり $u(x,1) = \frac{8}{\pi}\sum \frac{\sin((2k+1)x)}{(2k+1)^3}$ は $x(\pi-x)$ の標準的なフーリエ正弦級数と一致し、境界条件をすべて満たす。
